\documentclass[14pt,autoref,href,facsimile
%,fixint=false
%,times
]{disser}

\usepackage[a4paper,nohead,includefoot,mag=1000,
            margin=2cm,footskip=1cm]{geometry}
\usepackage[T2A]{fontenc}
\usepackage[utf8x]{inputenc}
\usepackage[english,russian]{babel}
\usepackage{tabularx}
\ifpdf\usepackage{epstopdf}\fi

% Поддержка нескольких списков литературы в одном документе
\usepackage{multibib}
% Создание команд для цитирования собственных работ диссертанта
% в отдельном разделе. В данном случае ссылка будет иметь вид \citemy{...}.
\newcites{my}{Список публикаций}

% Путь к файлам с иллюстрациями
\graphicspath{{figures/}}


\usepackage{amsmath}
\usepackage{amsmath,amssymb,amsthm}
\usepackage{multicol}
\usepackage{color}
%\usepackage{graphicx}
\usepackage{ucs}
\usepackage{pb-diagram}
\usepackage{enumerate}
\usepackage{appendix}
\usepackage{listings}
\usepackage{textcomp}


\definecolor{listinggray}{gray}{0.9}
\definecolor{lbcolor}{rgb}{0.9,0.9,0.9}
\lstset{
%	backgroundcolor=\color{lbcolor},
        language=Mathematica,
	tabsize=2,
	rulecolor=,
        basicstyle=\scriptsize,
        upquote=true,
%        aboveskip={1.5\baselineskip},
        columns=fixed,
        showstringspaces=false,
        extendedchars=true,
%        breaklines=true,
        prebreak = \raisebox{0ex}[0ex][0ex]{\ensuremath{\hookleftarrow}},
%        frame=single,
        showtabs=false,
        showspaces=false,
        showstringspaces=false,
        identifierstyle=\ttfamily,
        keywordstyle=\color[rgb]{0,0,1},
        commentstyle=\color[rgb]{0.133,0.545,0.133},
        stringstyle=\color[rgb]{0.627,0.126,0.941},
}

\newtheorem{statement}{Утверждение}
\newtheorem{theorem}{Теорема}
\newtheorem{axiom}{Аксиома}
\newtheorem{corollary}{Следствие}[chapter]
\newtheorem{lemma}{Лемма}
\newtheorem{mynote}{Замечание}[chapter]
%%\newtheorem{Def}{Definition}[section]
\newtheorem{Def}{Определение}[chapter]
\newtheorem{Cnj}[Def]{Гипотеза}
\newtheorem{Prop}{Свойство}
%%\newtheorem{example}{Example}[section]

\theoremstyle{definition}
\newtheorem{definition}{Определение}
\newtheorem{remark}{Замечание}[chapter]
\newtheorem{example}{Пример}[chapter]
\newtheorem{exercise}{Упражнение}
\newtheorem{conjecture}{Гипотеза}[chapter]

\newcommand{\go}{\stackrel{\circ }{\mathfrak{g}}}
\newcommand{\ao}{\stackrel{\circ }{\mathfrak{a}}}
\newcommand{\co}[1]{\stackrel{\circ }{#1}}
\newcommand{\pia}{\pi_{\mathfrak{a}}}
\newcommand{\piab}{\pi_{\mathfrak{a}_{\perp}}}
\newcommand{\gf}{\mathfrak{g}}
\newcommand{\af}{\mathfrak{a}}
\newcommand{\uf}{\mathfrak{u}}
\newcommand{\sfr}{\mathfrak{s}}
\newcommand{\aft}{\widetilde{\mathfrak{a}}}
\newcommand{\afb}{\mathfrak{a}_{\perp}}
\newcommand{\hf}{\mathfrak{h}}
\newcommand{\hfb}{\mathfrak{h}_{\perp}}
\newcommand{\pf}{\mathfrak{p}}

\newcommand{\gfh}{\hat{\mathfrak{g}}}
\newcommand{\afh}{\hat{\mathfrak{a}}}
\newcommand{\bff}{\mathfrak{b}}
\newcommand{\hfg}{\hf_{\gf}}


\begin{document}
% Включение файла с общим текстом диссертации и автореферата
% (текст титульного листа и характеристика работы).
% Общие поля титульного листа диссертации и автореферата
\institution{Санкт-Петербургский Государственный Университет}

\topic{Правила ветвления аффинных алгебр Ли\\ и приложения в моделях конформной теории поля}

\author{Назаров Антон Андреевич}

\specnum{01.04.02}
\spec{Теоретическая физика}

\sa{Ляховский Владимир Дмитриевич}
\sastatus{д.~ф.-м.~н., проф.}

\city{Санкт-Петербург}
\date{\number\year}

% Общие разделы автореферата и диссертации

\mkcommonsect{actuality}{Актуальность работы}{%

  Последние тридцать лет конформная теория поля в двух измерениях привлекает большое внимание исследователей. Эта теория используется для описания критического поведения в двумерных статистических системах. Благодаря наличию бесконечномерной алгебры симметрии двумерная конформная теория поля может быть сформулирована аксиоматически. Помимо математической красоты теория обладает огромной практической ценностью -- с ее использованием было получено большое количество результатов и численных предсказаний в изучении критического поведения в двумерных системах \cite{difrancesco1997cft,henkel1999conformal}. Методы двумерной конформной теории поля с успехом применяются также при изучении эффекта Кондо \cite{cox1998exotic,affleck1993exact} и дробного квантового эффекта Холла \cite{moore1991nonabelions}. 

Поиски строгого математического доказательства для предсказаний двумерной конформной теории поля \cite{cardy1992critical} в последние годы привели к большому количеству новых идей и результатов в дискретном комплексном анализе \cite{smirnov2001critical,duminil2011conformal,smirnov2010discrete}.

Теория представлений бесконечномерных алгебр Ли является важным инструментом изучения моделей конформной теории поля. Помимо алгебры Вирасоро, наличие которой обязательно в двумерной конформной теории поля, большую роль играют аффинные алгебры Ли. Изучение аффинных алгебр Ли было начато Виктором Кацем и Робертом Муди в 1960-х годах с попытки обобщения классификации простых конечномерных алгебр Ли на бесконечномерный случай \cite{kac1968simple,moody1968new}. Первоначально интерес к этим алгебрам был связан с модулярными свойствами характеров их модулей \cite{kac1984infinite,macdonald1971affine}. После возникновения двумерной конформной теории поля были предложены модели Весса-Зумино-Новикова-Виттена \cite{witten1984nab}, а затем и coset-модели \cite{Goddard198588}, в которых теория представлений аффинных алгебр Ли играет определяющую роль. 

Моделям Весса-Зумино-Новикова-Виттена, coset-моделям и теории представлений аффинных алгебр Ли  посвящены тысячи работ. Однако многие проблемы по-прежнему не имеют простых решений. Например, задача вычисления коэффициентов ветвления для представлений алгебр Ли стоит уже многие десятилетия. Она актуальна для различных физических приложений в coset-моделях конформной теории поля. При этом, в отличие от проблемы вычисления кратностей весов, для вычисления коэффициентов ветвления не существовало особенно эффективных алгоритмов. 
}

  
\mkcommonsect{objective}{Цели и задачи работы}{%
 Разработка рекуррентного подхода к функциям ветвления аффинных алгебр Ли, его связь с проблемами теории представлений и его приложения в моделях конформной теории поля.
}
 
\mkcommonsect{novelty}{Научная новизна и практическая значимость.}{%
  В диссертации впервые решены следующие задачи:
  \begin{itemize}%%[topsep=0pt, itemsep=-1ex]
  \item Получено эффективное рекуррентное соотношение для коэффициентов ветвления модулей аффинных и конечномерных алгебр Ли на модули не максимальных подалгебр.   Алгоритм вычисления коэффициентов ветвления реализован в пакете  {\bf Affine.m} для популярной системы компьютерной алгебры {\it Mathematica}.
  \item Установлена прямая связь инъективного сплинта и ветвлений. Доказано, что при определенных условиях  кратности весов вспомогательного модуля иньективного сплинта совпадают с коэффициентами ветвления в редукции на вложенную подалгебру. Наличие расщепления приводит к существенному упрощению при вычислении коэффициентов ветвления. 
%  \item Показано, что наличие расщепления приводит к существенному упрощению при вычислении коэффициентов ветвления.
  \item Исследована связь процедуры редукции с обобщенной резольвентой Бернштейна-Гельфанда-Гельфанда (БГГ).  Показано, что разложение сингулярного элемента определяет как коэффициенты ветвления, так и обобщенную БГГ-резольвенту, так как действие веера вложения на компоненты разложения порождает обобщенные модули Верма, которые образуют точную последовательность.
  \item Построена модель обобщенного стохастического процесса Шрамма-Лёвнера для систем с калибровочной инвариантностью, соответствующих coset-моделям конформной теории поля.
  \end{itemize}
Отметим, что  пакет {\bf Affine.m} может быть использован для решения задач теории представлений конечномерных и аффинных алгебр Ли, возникающих в различных областях физики, начиная от изучения атомных и молекулярных спектров и заканчивая конформной теорией поля и интегрируемыми системами.
}
 
%\mkcommonsect{value}{Практическая значимость}{%
%}

\mkcommonsect{results}{На защиту выносятся следующие результаты и положения:}{%
\begin{itemize}%%[topsep=0pt, itemsep=-1ex]
\item Получены новые рекуррентные соотношения на коэффициенты ветвления представлений аффинных алгебр Ли на представления произвольных редуктивных подалгебр, с использованием разложения сингулярных элементов
\item Установлено, что разложение сингулярного элемента определяет как коэффициенты ветвления, так и обобщенную БГГ-резольвенту, так как действие веера вложения на компоненты разложения порождает обобщенные модули Верма, которые образуют точную последовательность
\item  Доказано, что при определенных условиях  кратности весов вспомогательного модуля иньективного сплинта совпадают с коэффициентами ветвления в редукции на вложенную подалгебру. Наличие расщепления приводит к существенному упрощению при вычислении коэффициентов ветвления.
\item Показано, что условие для мартингала, определяющее классификацию операторов изменения граничных условий в наблюдаемых стохастического процесса Шрамма-Лёвнера, задает ограничения на структуру сингулярных элементов представлений аффинной алгебры Ли, порожденных граничными состояниями. Изучение структуры сингулярных элементов существенно упрощает поиск операторов смены граничных условий. Построена модель обобщенного стохастического процесса Шрамма-Лёвнера для систем с калибровочной инвариантностью, соответствующих coset-моделям конформной теории поля и показано, что такое обобщение совместно с coset-реализацией минимальных моделей.
%\item Продемонстрирована роль сингулярных элементов в построении мартингалов стохастического процесса Шрамма-Лёвнера, то есть проиллюстрировано применение алгебраических методов теории представлений аффинных алгебр Ли в изучении критического поведения в двумерных решеточных моделях
\item Разработан пакет программ {\bf Affine.m}, реализующий различные алгоритмы для вычислений в теории представлений конечномерных и аффинных алгебр Ли
\end{itemize}

}

\mkcommonsect{approbation}{Апробация работы}{%
  Материалы диссертации докладывались на трех международных конференциях, а также на семинарах кафедры физики высоких энергий и элементарных частиц СПбГУ, на семинарах в лаборатории имени П.Л. Чебышева математико-механического факультета СПбГУ, на семинаре лаборатории теоретической физики ОИЯИ (Дубна).

%  международном семинаре молодых ученых ``Workshop on Advanced Computer Simulation Methods''  27 - 29 апреля 2009 
%(Санкт-Петербург),  на международных конференциях:  ``Модели квантовой теории поля (MQFT-2010)'' 18-22 октября 2010 
%(Санкт-Петербург), ``Supersymmetries and Quantum Symmetries - 2011'', 18-23 июля 2011 (Дубна), ``Quantum Theory and 
%Symmetries (QTS-7)'', 7-13 августа 2011 (Прага). 
}

\mkcommonsect{pub}{Публикации.}{%
Материалы диссертации опубликованы в $10$ печатных работах, из них $5$ статей в
рецензируемых журналах~\citemy{2010arXiv1007.0318L,2011arXiv1102.1702L,2011arXiv1111.6787L,NazarovJETPletters,2011arXiv1107.4681N}, $5$ статей в
сборниках тезисов и трудов конференций \citemy{2011arXiv1112.4354N,2010LyakhovskyNazarovMQFT,Nazarov2008,NazarovACSM2009,2012arXiv1204.1855L}.
}

\mkcommonsect{contrib}{Личный вклад автора.}{%
  Все основные результаты и выносимые на защиту положения получены автором самостоятельно. Личный вклад автора в работы с соавтором составляет $50$ процетнов, в работы без соавторов -- $100$ процентов. 
}

\mkcommonsect{struct}{Структура и объем диссертации}{%
Диссертация состоит из введения и пяти глав, содержит 160 страниц и 30 рисунков. Список литературы включает 151 наименование. 

}

%%  \mkcommonsect{content}{Содержание работы}{%
%%  
%%  }

%%
%% End of file
%%% Local Variables: 
%%% mode: latex
%%% TeX-master: "thesis"
%%% End: 


\title{АВТОРЕФЕРАТ\\
диссертации на соискание ученой степени\\
кандидата физико-математических наук}

\maketitle

% Внутренняя сторона обложки
\noindent
\begin{center}
Работа выполнена на кафедре физики высоких энергий и элементарных частиц физического факультета Санкт-Петербургского государственного университета.
% \emph{название организации}.
\end{center}
\vskip1ex
\begin{tabularx}{\linewidth}{lp{1cm}X}
Научный руководитель:  & & \emph{доктор физико-математических наук}, \\
                       & & \emph{профессор}, \\
                       & & \emph{Ляховский Владимир Дмитриевич}
\\
Официальные оппоненты: & & \emph{доктор физико-математических наук}, \\
                       & & \emph{профессор}, \\
                       & & \emph{Кулиш Петр Петрович}\\
                       & & \emph{кандидат физико-математических наук}, \\
                       & & \emph{ученое звание}, \\
                       & & \emph{Мудров Андрей И}
\\
Ведущая организация:   & & \emph{Объединенный институт ядерных исследований}\\
\end{tabularx}

\vskip2ex\noindent
Защита состоится \datefield{} в \rule[0pt]{1cm}{0.5pt}\, часов
на заседании совета \emph{Д 212.232.24} по защите докторских и кандидатских диссертаций при \emph{Санкт-Петербургском государственном университете}, расположенном по адресу:
\emph{Санкт-Петербург, Средний пр. В.О., д. 41/43, ауд. 305}

\vskip1ex\noindent
С диссертацией можно ознакомиться в научной библиотеке
\emph{Санкт-Петербургского государственного университета}.

\vskip1ex\noindent
Автореферат разослан \datefield{}

\vskip2ex\noindent
%Отзывы и замечания по автореферату в двух экземплярах, заверенные
%печатью, просьба высылать по вышеуказанному адресу на имя ученого секретаря
%диссертационного совета.

\vfill\noindent
Ученый секретарь\\
диссертационного совета,\\
\emph{ученая степень}, \emph{ученое звание}%
\hfill
\makeatletter
% вставка файла, содержащего факсимиле ученого секретаря
\ifDis@facsimile
  \raisebox{-4pt}{\includegraphics[width=3cm]{sec-facsimile}}\hfill
\fi%
\makeatother%
\emph{фамилия и. о.}

\clearpage

\nsection{Общая характеристика работы}

% Актуальность работы
\actualitysection
\actualitytext

% Цель диссертационной работы
\objectivesection
\objectivetext

% Научная новизна
\noveltysection
\noveltytext

% Практическая значимость
%\valuesection
%\valuetext

% Результаты и положения, выносимые на защиту
\resultssection
\resultstext

% Апробация работы
\approbationsection
\approbationtext

% Публикации
\pubsection
\pubtext

% Личный вклад автора
%\contribsection
%\contribtext

% Структура и объем диссертации
\structsection
\structtext

% \contentsection
% \contenttext



\nsection{Содержание работы}


\textbf{Во Введении} обоснована актуальность диссертационной работы,
сформулирована цель и аргументирована научная новизна исследований, показана
практическая значимость полученных результатов, представлены выносимые на
защиту научные положения.

\textbf{Глава 1} носит обзорный характер. В ней мы даем аксиоматическую формулировку конформной теории поля, описываем модели Весса-Зумино-Новикова-Виттена и coset-модели. Затем мы демонстрируем роль аффинных алгебр в описании этих моделей и приводим основные понятия теории представлений, использующиеся в диссертации. Кроме того, мы обсуждаем конформную теорию поля на области с границей, так как она оказывается связана со стохастическим описанием решеточных моделей. 

\textbf{В главе 2} мы показываем, что структура сингулярного элемента определяет свойства модуля алгебры Ли, доказываем лемму о разложении сингулярного элемента и выводим основное рекуррунтное соотношение на коэффициенты ветвления. Основные результаты данной главы опубликованы в работе \citemy{2010arXiv1007.0318L}. 

Формула Вейля-Каца для формальных характеров интегрируемых модулей старшего веса конечномерных и аффинных алгебр Ли имеет вид
\begin{equation}
  \label{eq:1}
  \mathrm{ch} V^{(\mu)} = \frac{\Psi^{(\mu)}}{R},
\end{equation}
где $\Psi^{(\mu)}$ -- сингулярный элемент модуля, а $R=\prod_{\alpha\in \Delta^+}(1-e^{-\alpha})^{\mathrm{mult}(\alpha)}$ -- знаменатель Вейля. Сингулярный элемент определяется набором сингулярных весов модуля и имеет разный вид для разных типов модулей старшего веса. Например, $\Psi^{(\mu)}=\sum_{w\in W} \epsilon(w) e^{w(\mu+\rho)-\rho}$ для неприводимых модулей. Заметим, что знаменатель Вейля является универсальным объектом, характеризующим корневую систему алгебры Ли, а свойства модуля определяются сингулярным элементом.

Процедура редукции состоит в разложении модуля алгебры Ли $\gf$ в сумму модулей некоторой подалгебры $\af$
\begin{equation}
  \label{eq:2}
  L_{\gf\downarrow \af}^{\mu }=\bigoplus
\limits_{\nu \in P_{\af}^{+}}b_{\nu }^{\left( \mu \right) }L_{\af}^{\nu }.
\end{equation}
Используя оператор проекции  $\pi_{\af}$ (на весовое пространство $\hf_{\af}^*$), перепишем это разложение для формальных характеров:
\begin{equation}
\label{branching1}
 \pi _{\af}\circ ch\left( L^{\mu }\right)
 =\sum_{\nu \in P_{\af}^{+}}b_{\nu }^{(\mu)}ch\left( L_{\af}^{\nu }\right) .
\end{equation}
Нас интересуют коэффициенты ветвления $b^{(\mu)}_{\nu}$.

Для любой алгебры $\gf$ и подалгебры $\af\subset \gf$ можно построить  подалгебру $\afb$ такую, что 
\begin{eqnarray}
\label{delta-a-ort}
\Delta _{\af_{\perp }} &:&=\left\{ \beta \in \Delta _{\gf}|
\forall h \in \hf_{\af};  \beta\left(h \right)=0  \right\} , \\
\Delta _{\af_{\perp }}^{+} &:&=\left\{ \beta ^{+}\in \Delta _{\gf%
}^{+}|\forall h \in \hf_{\af};  \beta^{+}\left(h \right)=0  \right\} .
\end{eqnarray}
Обозначим через $W_{\afb}$ подгруппу группы Вейля $W$, порожденную отражениями $w _{\beta }$, соответствующими корням $\beta \in \Delta _{\afb}^{+}$ . Подсистема  $\Delta _{\af_{\perp }}$ определяет подалгебру $\af_{\perp }$ с подалгеброй Картана $\hf_{\afb}$. 

Мы доказываем следующую лемму о разложении сингулярного элемента:
\begin{lemma}
\label{lemma}
Пусть $\left( \af,\afb \right)$ -- ортогональная пара редуктивных подалгебр $\gf$ и  $\widetilde{\afb}=\afb\oplus \hf_{\perp }$, $\widetilde{\af}=\af\oplus\hf_{\perp }$ ,

$L^{\mu }$ -- модуль старшего веса с сингулярным элементом $\Psi ^{\left(\mu \right)}$ ,

$R_{\af_{\perp }}$ -- знаменатель Вейля для подалгебры $\af_{\perp }$.

Тогда элемент  $\Psi ^{\left( \mu \right) }_{\left(  \af, \afb \right)}=\pi _{\af}\left( \frac{\Psi _{\gf}^{\mu }}{R_{\af_{\perp }}}\right) $ можно разложить в сумму по  $u\in U$ (см. (\ref{U-def})) сингулярных весов $e^{\mu _{\af}\left( u\right) }$ с коэффициентами $\epsilon (u)\mathrm{\dim}\left( L_{\widetilde{\afb}}^{\mu _{\widetilde{\afb}}\left( u\right) }\right) $:
\begin{equation}
\Psi ^{\left( \mu \right) }_{\left(  \af, \afb \right)}=\quad \pi _{\af}\left( \frac{\Psi^{\mu }}{R_{\af%
_{\perp }}}\right) =\sum_{u\in U}\;\epsilon (u)\mathrm{\dim }
\left( L_{\widetilde{\af_{\perp }}}^{\mu _{%
\widetilde{\af_{\perp }}}\left( u\right) }\right) e^{\mu _{\af}\left( u \right) }.
\end{equation}
\end{lemma}

Теперь мы можем ввести ``веер вложения'', который необходим для формулировки рекуррентных соотношений:
\begin{definition}
\label{fan-definition} Рассмотрим произведение
\begin{equation}
\prod_{\alpha \in \Delta ^{+}\setminus \Delta _{\afb }^{+}}\left( 1-e^{-\pi
_{\af}\alpha }\right) ^{\mathrm{mult}(\alpha )-\mathrm{mult}_{\af%
}(\pi _{\af}\alpha )}=-\sum_{\gamma \in P_{\af}}s(\gamma
)e^{-\gamma }  \label{eq:142}
\end{equation}
и носитель $\Phi _{\af\subset \gf}\subset P_{\af}$ функции $s(\gamma )=\det \left( \gamma \right) $ :
\begin{equation}
\Phi _{\af\subset \gf}=\left\{ \gamma \in P_{\af}|s(\gamma
)\neq 0\right\}   \label{eq:37}
\end{equation}
Упорядочение корней в  $\co{\Delta _{\af}}$ индуцирует естественное упорядочение весов в $P_{\af}$. Обозначим через $\gamma_{0}$ наименьший вектор $\Phi _{\af\subset \gf}$. Множество
\begin{equation}
\Gamma _{\af\rightarrow \gf}=\left\{ \xi -\gamma _{0}|\xi \in \Phi _{%
\af\subset \gf}\right\} \setminus \left\{ 0\right\}
\label{fan-defined}
\end{equation}
называется  \textit{веером вложения}.
\end{definition}
Веер вложения универсален и зависит только от вложения $\af\to\gf$ и не зависит от модуля $L^{(\mu)}$.

Числитель описывает ветвление в терминах сингулярных элементов и, как элемент  $\mathcal{E}\left( \gf \right)$, допускает разложение в виде:
\begin{equation}
  \label{eq:21}
  \sum_{\nu \in \bar{C_{\mathfrak{a}}}}b_{\nu }^{\left( \mu \right) }\Psi _{\left( \frak{%
        a}\right) }^{\left( \nu \right) }=\sum_{\lambda \in P_{\af}}k_{\lambda
  }^{\left( \mu \right) }e^{\lambda }.
\end{equation}
Здесь $k_{\lambda}^{\left( \mu \right) }$  -- целочисленные коэффициенты, их знаки зависят от длины элементов группы Вейля в 
$\Psi _{\left( \af\right) }^{\left( \nu \right) }$. Важно заметить, что $k_{\lambda}^{\left( \mu \right) }$ совпадают с коэффициентами ветвления для всех весов $\nu$ в главной камере Вейля:
\begin{equation}
  \label{eq:119}
  b^{(\mu)}_{\nu}=k^{(\mu)}_{\nu} \; \mbox{для} \; \nu\in \bar{C}_{\mathfrak{a}}.
\end{equation}
Мы называем коэффициенты $k_{\lambda}$ сингулярными коэффициентами ветвления
(см. также \cite{ilyin812pbc}).

Теперь мы можем сформулировать основную в этой главе теорему, которая позволит нам рекуррентно вычислять коэффициенты ветвления.

\begin{theorem}
  Для сингулярных коэффициентов ветвления $k^{(\mu)}_{\nu}$ (\ref{eq:21}) выполняется соотношение
  \begin{equation}
    \label{recurrent-relation}
    \begin{array}{c}
      k_{\xi }^{\left( \mu \right) }=-\frac{1}{s\left( \gamma _{0}\right) }\left(
        \sum_{u\in U} \epsilon(u)\;
        \dim \left( L_{\widetilde{\af_{\perp }}}^{\mu
        _{\widetilde{\af_{\perp }}}\left( u\right) }\right)
        \delta_{\xi-\gamma_0,\pi_{\af}(u(\mu+\rho)-\rho)}+ \right.\\
      \left.
        +\sum_{\gamma \in
          \Gamma _{\af \rightarrow \gf}}s\left( \gamma +\gamma _{0}\right) k_{\xi
          +\gamma }^{\left( \mu \right) }\right).
    \end{array}
  \end{equation}
\end{theorem}

Далее мы анализируем пары $(\af,\afb)$ для простых алгебр Ли. Оказывается, что для ``ортогональной пары'' $(\af,\afb)$, вообще говоря, $\af\oplus\afb\not\subset\gf $. В частности, для серии простых конечномерных алгебр $B_n$ существуют ``ортогональные пары'' подалгебр $(B_k,B_{n-k})$.

На основании рекуррентного соотношения (\ref{recurrent-relation}) сформулирован следующий алгоритм вычисления коэффициентов ветвления:

\begin{enumerate}[(i)]
\item Построить корневую систему $\Delta _{\af}$ для вложения $\af\rightarrow \gf$.

\item Выбрать все положительные корни $\alpha \in \Delta ^{+}$, ортогональные к  $\af$, то есть сформировать множество $\Delta_{\afb }^{+}$.

\item Построить множество $\Gamma _{\af\rightarrow \gf}$. Соотношение  (\ref{eq:142}) определяет знаковую функцию $s(\gamma)$ и множество $\Phi_{\af\subset \gf}$. Вычитание наименьшего веса $\gamma_0$ дает веер вложения (\ref{fan-defined}):
 $\Gamma _{\af\rightarrow \gf}=\left\{ \xi -\gamma _{0}|\xi \in \Phi _{\af\subset \gf}\right\} \setminus \left\{ 0\right\}$.

\item Построить множество $\widehat{\Psi ^{(\mu )}}=\left\{ w (\mu +\rho
)-\rho ;\;w \in W\right\} $ сингулярных весов для  $\gf$-модуля $L^{(\mu )}$.

\item Выбрать веса $\left\{ \mu _{\widetilde{\afb}}\left( w\right) =\pi _{\widetilde{\afb}}\left[ w(\mu +\rho
)-\rho \right] -\mathcal{D}_{\afb}\in \overline{C_{\widetilde{\afb}}}\right\} $. Так как множество  $\Delta_{\afb }^{+}$ фиксировано, легко осуществить проверку принадлежности веса $\mu _{\widetilde{\afb}}\left( w\right) $ к главной камере Вейля $\overline{C_{\widetilde{\afb}}}$ (достаточно вычислить скалярные произведения с фундаментальными весами $\afb^{+}$).

\item Для весов $\mu _{\widetilde{\afb}}\left( w\right) $ вычислить размерности соответствующих модулей, $\mathrm{\dim }\left(L_{\widetilde{\afb}}^{\mu _{\widetilde{\afb}}\left( u\right) }\right) $, используя формулу Вейля для размерностей. Затем построить сингулярный элемент $\Psi ^{\left( \mu \right) }_{\left(  \af, \afb \right)}$.

\item Вычислить сингулярные коэффициенты ветвления используя рекуррентное соотношение (\ref{recurrent-relation}) и выбрать среди них те, которые соответствуют весам в главной камере Вейля $\overline{C_{\af}}$.
\end{enumerate}

Остальные разделы главы 2 содержат примеры вычислений с использованием предложенного алгоритма, а также описание роли функций ветвления в формулировке конформной теории поля на торе и в coset-моделях конформной теории поля. 


В следующей \textbf{главе 3} мы используем  разложение сингулярного элемента, чтобы показать связь ветвления с (обобщенной) резольвентой Бернштейна-Гельфанда-Гельфанда.   Результаты третьей главы опубликованы в работах \citemy{2011arXiv1102.1702L,2010LyakhovskyNazarovMQFT}.

%%  3 Коэффициенты ветвления и обобщенная резольвента Бернштейна-Гельфанда-Гельфанда
%%    3.1 Ортогональная подалгебра и сингулярные элементы
%%      3.1.1 Разложение сингулярного элемента.
%%      3.1.2 Фомулы Вейля-Верма.
%%    3.2 БГГ резольвента и ветвление
%%    3.3 Заключение

Оказывается, что для полупростой конечномерной алгебры $\gf$  и полупростой конечномерной подалгебры $\af$ алгебра $\afb$, определенная в формуле (\ref{branching1}), является регулярной. Тогда отношение знаменателей Вейля порождает параболические модули Верма.



%%     4 Сплинты корневых систем и функции ветвления
%%       4.1 Вложения и сплинты
%%       4.2 Как стебли определяют функции кратности
%%       4.3 Примеры
%%       4.4 Сплинты и соотношения для аффинных алгебр Ли
%%       4.5 Заключение
%%       4.6 Приложение
%%     5 Практические приложения
%%       5.1 Coset-модели и критическое поведение
%%         5.1.1 Введение
%%         5.1.2 Эволюция Шрамма-Левнера
%%         5.1.3 Мартингалы эволюции Шрамма-Левнера и корреляционные функции в конформной теории поля
%%         5.1.4 Модели Весса-Зумино-Новикова-Виттена и эволюция Шрамма-Левнера
%%         5.1.5 Coset-модели и эволюция Шрамма-Левнера
%%           5.1.5.1 Примеры
%%         5.1.6 Coset-модели и интегрируемые деформации конформной теории поля
%%         5.1.7 О связи мартингалов coset-моделей с классификацией операторов смены граничного условия
%%       5.2 Вычислительные методы для теории представлений аффинных алгебр Ли
%%         5.2.1 Структуры данных
%%           5.2.1.1 Веса
%%           5.2.1.2 Корневые системы
%%           5.2.1.3 Формальные элементы
%%           5.2.1.4 Модули
%%         5.2.2 Основные алгоритмы
%%       5.3 Примеры вычислений
%%         5.3.1 Разложение тензорных произведений модулей конечномерных алгебр Ли на неприводимые
%%         5.3.2 Ветвления и параболические модули Верма
%%         5.3.3 Струнные функции аффинных алгебр Ли и модели конформной теории поля
%%         5.3.4 Функции ветвления и coset-модели конформной теории поля
%%       5.4 Заключение
%%  


\textbf{Глава 4} посвящена сплинтам -- расщеплением корневой системы алгебры Ли в объединение образов корневых систем двух алгебр, не обязательно являющихся подалгебрами данной алгебры. Если одна из алгебр является подалгеброй, то сплинт приводит к резкому упрощению в вычислении коэффициентов ветвления -- они совпадают с кратностями весов в модуле другой алгебры. Основная часть главы посвящена доказательству этого факта. Кроме того, сплинт корневой системы простой конечномерной алгебры Ли приводит к возникновению новых соотношений на струнные функции и функции ветвления соответствующего аффинного расширения. Эти соотношения обсуждаются в разделе 4.4.
Данные результаты опубликованы в статьях \citemy{2011arXiv1111.6787L,2012arXiv1204.1855L}.

Заключительная \textbf{глава 5} посвящена практическим приложениям результатов диссертации. В разделе 5.1 мы  описываем применение алгебраических методов к проблеме поиска соответствия между квантовополевым и решеточным описанием критического поведения. Эти результаты были опубликованы нами в работах \citemy{NazarovJETPletters,2011arXiv1112.4354N}. Раздел 5.2 представляет собой описание пакета {\bf Affine.m}, предназначенного для вычислений в теории представлений аффинных и конечномерных алгебр Ли и реализованного с использованием методов диссертации. Вычислительным методам посвящены наши работы \citemy{2011arXiv1107.4681N,NazarovACSM2009,Nazarov2008}.

%%  
%%  \textbf{Во Введении} обоснована актуальность диссертационной работы,
%%  сформулирована цель и аргументирована научная новизна исследований, показана
%%  практическая значимость полученных результатов, представлены выносимые на
%%  защиту научные положения.
%%  
%%  \textbf{В первой главе} ...
%%  
%%  Содержание первой главы.
%%  
%%  Результаты первой главы опубликованы в
%%  работе~\cite{2010arXiv1007.0318L}
%%  
%%  \textbf{Во второй главе} ...
%%  
%%  Содержание второй главы.
%%  
%%  Результаты второй главы опубликованы в
%%  работе~\citemy{Petrov_2001_Journal_23_12321}.
%%  
%%  \textbf{В третьей главе} ...
%%  
%%  Содержание третьей главы.
%%  
%%  Результаты третьей главы опубликованы в
%%  работе~\citemy{Sidorov_2002_Journal_32_1531}.
%%  
%%  \textbf{В Заключении}
%%  
% ----------------------------------------------------------------
\renewcommand\bibsection{\nsection{Список публикаций}}

% Префикс номеров ссылок на работы соискателя
\def\BibPrefix{A}
\bibliographystylemy{disser}
\bibliographymy{bibliography}

\renewcommand\bibsection{\nsection{Цитированная литература}}

\def\BibPrefix{}
\bibliographystyle{disser}
\bibliography{bibliography}
% ----------------------------------------------------------------

\end{document}
